\chapter{Constipation}

\section*{Definition}
Constipation commonly involves fewer than three spontaneous bowel movements per week, hard stools, straining, a feeling of incomplete evacuation, or the need for manual manoeuvres. Frequency alone does not establish the diagnosis.

\section*{When to See a Doctor}
See your doctor if:
\begin{itemize}
  \item You have rectal bleeding, a change in your usual bowel pattern, or a family history of colon cancer
\end{itemize}

\section*{How It Develops}
It can result from slow colonic transit, difficulty coordinating pelvic-floor muscles, diet or activity changes, medicines, endocrine or neurologic disease, or a disorder such as constipation-predominant IBS.

\section*{Causes}
\begin{itemize}
  \item Low fibre intake, inadequate fluids, or low physical activity
  \item Opioid medication use
  \item Hypothyroidism
  \item Parkinson's disease
  \item IBS-C constipation predominant
\end{itemize}

\section*{Related Symptoms}
\begin{itemize}
  \item Bloating, straining, hard stools, a sensation of blockage, or incomplete emptying
\end{itemize}

\section*{Medical Evaluation}
\begin{itemize}
  \item History assesses stool frequency and form, straining, medicines, diet, alarm symptoms, and possible pelvic-floor dysfunction.
  \item Examination may include abdominal and rectal examination, including assessment for fecal impaction when indicated.
  \item Blood tests, imaging, anorectal testing, or colonoscopy are selected for alarm features, suspected secondary causes, or persistent symptoms.
\end{itemize}

\section*{Treatment Options}
\begin{itemize}
  \item Management may include fibre, osmotic or stimulant laxatives, prescription secretagogues, or biofeedback for defecatory disorders.
  \item Opioid-induced constipation, hypothyroidism, obstruction, inflammatory disease, and other secondary causes require cause-specific treatment.
\end{itemize}

\section*{Self-Care and Home Management}
\begin{itemize}
  \item Increase fibre gradually through food and maintain adequate fluids unless a clinician has restricted your intake.
  \item Regular physical activity and responding to the urge to defecate may help; a footstool can improve toilet posture.
  \item Bulk fibre or an osmotic laxative such as polyethylene glycol may help when appropriate; follow the product label or pharmacist advice.
  \item Seek prompt care for vomiting, severe or worsening pain, marked distension, inability to pass gas, rectal bleeding, or unexplained weight loss.
\end{itemize}

\section*{References}
\begin{thebibliography}{99}
\bibitem{constipation:ref1} Bharucha AE, Pemberton JH, Locke GR III. ``American Gastroenterological Association technical review on constipation.'' Gastroenterology. 2013;144(1):218-238.
\bibitem{constipation:ref2} Camilleri M, Ford AC, Mawe GM, et al. ``Chronic constipation.'' Nat Rev Dis Primers. 2017;3:17095.
\bibitem{constipation:ref3} Lacy BE, Mearin F, Chang L, et al. ``Bowel disorders.'' Gastroenterology. 2016;150(6):1393-1407.
\bibitem{constipation:ref4} Rao SSC, Rattanakovit K, Patcharatrakul T. ``Diagnosis and treatment of chronic constipation.'' J Neurogastroenterol Motil. 2016;22(2):201-211.
\bibitem{constipation:ref5} Lindberg G, Hamid SS, Malfertheiner P, et al. ``World Gastroenterology Organisation global guideline: constipation.'' J Clin Gastroenterol. 2011;45(6):483-487.
\bibitem{constipation:ref6} Bharucha AE. ``Management of chronic constipation in adults.'' UpToDate. Wolters Kluwer; 2025.
\end{thebibliography}
